\section{The \texttt{Filter} McXtrace Component}
Release: McXtrace 1.1

Block of an attenuating material

\subsection*{Identification}
\begin{itemize}
  \item \textbf{Author:} Erik Knudsen
  \item \textbf{Origin:} DTU Physics
  \item \textbf{Date:} Jan 24, 2011
\end{itemize}

\subsection*{Description}
A chunk of attenuating material. Attenuation is computed through the effective length travelled within the material. No scattering is modelled at present.

Filter shape may be a cylinder, a sphere, a box or any other shape.

\begin{verbatim}
box/plate:       xwidth x yheight x zdepth
cylinder:        radius x yheight (along Y axis)
sphere:          radius
any shape:       geometry=OFF/PLY_file
\end{verbatim}

Example: Filter(material\_datafile="Ge.txt", geometry="wire.ply",xwidth=0.02,yheight=0,zdepth=0) Example: Filter(material\_datafile="Ge.txt",xwidth=0.02,yheight=0.02, zdepth=1e-4) Example: Filter(material\_datafile="Ge.txt",radius=1e-4,yheight=0.02) Example: Filter(material\_datafile="Ge.txt",radius=1e-3, refraction=1)

\subsection*{Input parameters}
Parameters in \textbf{boldface} are required; the others are optional.

\begin{longtable}{p{0.22\textwidth}p{0.12\textwidth}p{0.46\textwidth}p{0.14\textwidth}}
\toprule
\textbf{Name} & \textbf{Unit} & \textbf{Description} & \textbf{Default} \\
\midrule
\endhead
refraction & 0/1 & If nonzero, refraction is enabled. (Only functional for basic geometries, does not yet work for OFF) & 0 \\
fixed\_delta & 0/1 & Use a fixed delta to compute refraction - useful for debugging. & 0 \\
material\_datafile & str & File where the material parameters for the filter may be found. Format is similar to what may be found off the NIST website. [Be.txt] & "Be.txt" \\
geometry & str & File containing the polygon definition of a general shape object. When xwidth is also given, the object is rescaled accordingly (OFF/PLY) & 0 \\
xwidth & m & Width of block. & 0 \\
yheight & m & Height of block. & 0 \\
zdepth & m & Thickness of block. & 0 \\
radius & m & Radius of cylinder or sphere. & 0 \\
mu\_col & idx & Column index to pick up absorption length mu, counted from 0. Use with non-standard input file, e.g. 2-3 column files. -1 means attempt autodetect. & -1 \\
\bottomrule
\end{longtable}

\subsection*{Links}
\begin{itemize}
  \item Component source code found in file \texttt{Filter.comp}.
  \item \htmladdnormallink{Meshlab}{https://www.meshlab.net/} - viewer for OFF files
  \item \htmladdnormallink{Geomview and Object File Format (OFF)}{http://www.geomview.org} - OFF file definition and examples
  \item \htmladdnormallink{jroff.jar - Java version of Geomview (display only)}{http://www.holmes3d.net/graphics/roffview/}
  \item \htmladdnormallink{Qhull}{http://qhull.org} - for calculating a convex hull from points.
  \item \htmladdnormallink{Powercrust}{https://www.cs.ucdavis.edu/\textasciitilde{}amenta/powercrust.html} - for reconstructing a solid geometry from a point cloud.
  \item Most material datafile inpus may obtained from \htmladdnormallink{NIST FFast}{https://physics.nist.gov/PhysRefData/FFast/form.html}
\end{itemize}
\IfFileExists{optics/Filter_static.tex}{\input{optics/Filter_static.tex}}{}